%==============================================================================
\section{Conjunto de Dados}
\label{sec:dataset}

% --- corresponde a 04_Dataset.md ---

O substrato primário é uma única campanha de envelhecimento acelerado
controlada por PID (\texttt{Tentative\_\allowbreak{}1\_\allowbreak{}20260524\_\allowbreak{}203756.csv}),
iniciada em 24/05/2026, amostrada a 1~Hz por \valDurationHours~h
(\valRawSamples{} linhas brutas, das quais \valDropouts{} são
\emph{dropouts} de queda de enlace). Após excluir o transiente de
estabilização térmica de 1.0~h, restam \valPostSoakSamples{} amostras. O DUT
é mantido em um regime permanente de alto estresse
($\Tdie\approx\valDutTempMean$~$^\circ$C, $\Vcore\approx\valDutVoltMean$~V)
sob controle PID, com um modelo de planta FOPDT identificado
$G(s)=1.56\,e^{-150.6 s}/(1307.2 s+1)$.

A \Cref{tab:dataset} resume os canais relevantes ao sensor e as relações
medidas. O \emph{slack} ocupa \valSlackLevels{} níveis inteiros ao longo da
faixa de \valSlackRange{} contagens, resultando em uma entropia discreta de
$\Ent{\slk}=\valHs$~bits. A temperatura é o principal fator reversível
($\operatorname{corr}(\slk,\Tdie)=\valCorrST$), com um acoplamento de tensão
mais fraco ($\valCorrSV$) e uma deriva de envelhecimento lenta e genuína
($\operatorname{corr}(\slk,\tm)=\valCorrSt$; inclinação bruta por mínimos
quadrados \valSlopeRaw{} contagens/h). Uma regressão linear do \emph{slack}
em $(\Tdie,\Vcore)$ explica $R^2=\valRsqLinear$ da variância --- um valor
que, como a física subjacente é não-linear, motivou a hipótese central
deste trabalho de que $R^2$ subestima o verdadeiro acoplamento ambiental,
testada --- e não confirmada, ao menos não como não-linearidade
instantânea --- na análise de informação mútua da \cref{sec:results}.

% Tabela-resumo do conjunto de dados (números de generated/values.tex).
\begin{table}[t]
  \centering
  \caption{Canais relevantes ao sensor e relações medidas na janela
           pós-estabilização ($n=\valPostSoakSamples$).}
  \label{tab:dataset}
  \begin{tabular}{lll}
    \toprule
    Grandeza & Valor & Leitura \\
    \midrule
    Duração              & \valDurationHours~h        & campanha PID única \\
    Amostras pós-soak     & \valPostSoakSamples        & após 1\,h de soak \\
    $\Tdie$ (média)      & \valDutTempMean~$^\circ$C & regime de alto estresse \\
    $\Vcore$ (média)     & \valDutVoltMean~V  & mantida pela PSU \\
    $\slk$ (média / desvio) & \valSlackMean\,/\,\valSlackStd & contagens \\
    $\slk$ níveis / faixa & \valSlackLevels\,/\,\valSlackRange & estados inteiros ocupados \\
    \midrule
    $\operatorname{corr}(\slk,\Tdie)$ & \valCorrST & temperatura domina \\
    $\operatorname{corr}(\slk,\Vcore)$ & \valCorrSV & fraco, sinal correto \\
    $\operatorname{corr}(\slk,\tm)$ & \valCorrSt & deriva lenta de envelhecimento \\
    inclinação bruta $d\slk/d\tm$ & \valSlopeRaw~cnt/h & confundida com PVT \\
    $R^2$ linear $(\slk\mid\Tdie,\Vcore)$ & \valRsqLinear & motivação da análise de MI (\cref{sec:res-mi}) \\
    $\Ent{\slk}$ (discreta) & \valHs~bits & orçamento total de saída \\
    \bottomrule
  \end{tabular}
\end{table}


Os riscos dos dados carregados para o \emph{pipeline} são: forte correlação
serial ($\neff\ll n$, com $\neff=\valNeff$ aqui), quantização inteira de
$\slk$, e um fator de confusão tendência/ambiente que exige remoção de
tendência antes da MI.

Para a comparação de bancadas da \cref{sec:res-kl}, são usadas duas
execuções \emph{estabilizadas} de $\sim$24~h: uma bancada SBCCI não-PID
(ruído de \emph{slack} $\valSlackStdSBCCI$ contagens) e a referência
controlada por PID ($\valSlackStdPID$ contagens). Ambas foram adquiridas em
um \textbf{dispositivo de teste dedicado, distinto do dispositivo de
envelhecimento de 214~h}; as duas bancadas compartilham esse dispositivo de
teste, de modo que sua divergência isola a diferença de bancada/ambiente,
mas seus níveis absolutos de \emph{slack} não são comparáveis aos do
dispositivo de envelhecimento. Seus pontos de operação distintos são
tratados comparando a flutuação centrada do \emph{slack}, de modo que a
divergência reflete o caráter do ruído, e não o ponto de ajuste.
